\documentclass[../main.tex]{subfiles}
\begin{document}
\subsection{S9 and S10: FOMC drift and turn-of-the-month}\label{sec:s9}

\paragraph{Economic intuition and literature.}
US equities may earn much of their excess return in the 24 hours before scheduled FOMC announcements \citep{lucca2015} and around month turns \citep{ariel1987,lakonishok1988}. The month-turn effect is consistent with payment-flow timing.

\paragraph{Specification.}
\begin{itemize}
\item \textbf{S9 (FOMC):} buy SPY at the close of the session before the announcement and sell at the close of announcement day. The data has no 14:00 price, so the post-announcement reaction is included (a known dilution). The portfolio is 100\% of equity over 24 events.
\item \textbf{S10 (turn of the month):} buy SPY at the close of the second-to-last session of the month and sell at the close of the third session of the next month. The portfolio is 100\% of equity.
\item \emph{Execution tier:} T1 official closes for both. Both are confirmatory but underpowered.
\end{itemize}

\paragraph{Implementation.}
The engine is \file{wave3/cpp/s_calendar.cpp}, with official-close fills. S9 has 24 events in total, so it cannot detect anything smaller than a very large effect.

\paragraph{Out-of-sample testing and results.}
Both strategies are positive in development (S9 $+2.3\%$ per year, Sharpe 0.53; S10 $+3.0\%$, Sharpe 0.36) and negative in 2024 (Table~\ref{tab:s9}). Both are rejected.

\begin{table}[tbp]
\centering\small
\caption{S9 and S10: calendar strategies on SPY.}\label{tab:s9}
\begin{tabular}{@{}llrrrr@{}}
\toprule
Strategy & Sample & Ann.\ excess & Sharpe & NW $t$ & \\
\midrule
S9 & Development & $+$2.3\% & 0.53 & 0.74 & \\
S9 & Test (2024) & $-$0.9\% & $-$0.25 & & \\
S10 & Development & $+$3.0\% & 0.36 & 0.55 & \\
S10 & Test (2024) & $-$3.6\% & $-$0.61 & & \\
\bottomrule
\end{tabular}
\end{table}
\ifSubfilesClassLoaded{\section*{References}% Bibliography.
\bigskip\noindent\emph{The reference details below should be checked against the published versions before citing.}\par\medskip
\begin{thebibliography}{99}
\setlength{\itemsep}{2pt}
\small
\bibitem[Bailey and L{\'o}pez de Prado(2014)]{bailey2014} Bailey, D. H., \& L{\'o}pez de Prado, M. (2014). The deflated Sharpe ratio. \emph{Journal of Portfolio Management}, 40(5).
\bibitem[Barber et~al.(2013)]{barber2013} Barber, B. M., De George, E. T., Lehavy, R., \& Trueman, B. (2013). The earnings announcement premium around the globe. \emph{Journal of Financial Economics}, 108(1).
\bibitem[Brandt et~al.(2008)]{brandt2008} Brandt, M. W., Kishore, R., Santa-Clara, P., \& Venkatachalam, M. (2008). Earnings announcements are full of surprises. Working paper.
\bibitem[Chan et~al.(1996)]{chan1996} Chan, L. K. C., Jegadeesh, N., \& Lakonishok, J. (1996). Momentum strategies. \emph{Journal of Finance}, 51(5).
\bibitem[Cremers and Weinbaum(2010)]{cremers2010} Cremers, M., \& Weinbaum, D. (2010). Deviations from put-call parity and stock return predictability. \emph{Journal of Financial and Quantitative Analysis}, 45(2).
\bibitem[DeMiguel et~al.(2009)]{demiguel2009} DeMiguel, V., Garlappi, L., \& Uppal, R. (2009). Optimal versus naive diversification: How inefficient is the 1/N portfolio strategy? \emph{Review of Financial Studies}, 22(5).
\bibitem[Easley et~al.(1998)]{easley1998} Easley, D., O'Hara, M., \& Srinivas, P. S. (1998). Option volume and stock prices: Evidence on where informed traders trade. \emph{Journal of Finance}, 53(2).
\bibitem[Fama and MacBeth(1973)]{famamacbeth1973} Fama, E. F., \& MacBeth, J. D. (1973). Risk, return, and equilibrium: Empirical tests. \emph{Journal of Political Economy}, 81(3).
\bibitem[Frazzini and Lamont(2007)]{frazzini2007} Frazzini, A., \& Lamont, O. A. (2007). The earnings announcement premium and trading volume. NBER Working Paper 13090.
\bibitem[Johnson and So(2012)]{johnsonso2012} Johnson, T. L., \& So, E. C. (2012). The option to stock volume ratio and future returns. \emph{Journal of Financial Economics}, 106(2).
\bibitem[Jorion(1986)]{jorion1986} Jorion, P. (1986). Bayes-Stein estimation for portfolio analysis. \emph{Journal of Financial and Quantitative Analysis}, 21(3).
\bibitem[Ledoit and Wolf(2004)]{ledoit2004} Ledoit, O., \& Wolf, M. (2004). A well-conditioned estimator for large-dimensional covariance matrices. \emph{Journal of Multivariate Analysis}, 88(2).
\bibitem[Markowitz(1952)]{markowitz1952} Markowitz, H. (1952). Portfolio selection. \emph{Journal of Finance}, 7(1).
\bibitem[Newey and West(1987)]{neweywest1987} Newey, W. K., \& West, K. D. (1987). A simple, positive semi-definite, heteroskedasticity and autocorrelation consistent covariance matrix. \emph{Econometrica}, 55(3).
\bibitem[Pan and Poteshman(2006)]{pan2006} Pan, J., \& Poteshman, A. M. (2006). The information in option volume for future stock prices. \emph{Review of Financial Studies}, 19(3).
\bibitem[Xing et~al.(2010)]{xing2010} Xing, Y., Zhang, X., \& Zhao, R. (2010). What does the individual option volatility smirk tell us about future equity returns? \emph{Journal of Financial and Quantitative Analysis}, 45(3).
\end{thebibliography}
}{}
\end{document}
